% 01-metodologia-pt.tex — Métodos
\section{Métodos}
\label{sec:methods}

\subsection{Sistema sob teste}
\label{sec:sistema}
Todos os experimentos foram executados dentro do OpenCode Ecosystem Core, um
ecossistema multiagente com 206 agentes de catálogo, um Blackboard A2A, um
bus metacognitivo (MetaBus), um Trust Engine e uma política estrita de
SDD/TDD. O orquestrador (MarceloClaro) executa o ciclo
Perceber$\rightarrow$Especificar$\rightarrow$Delegar$\rightarrow$Executar
$\rightarrow$Verificar$\rightarrow$Refletir. LLMs externos são invocados via
CLI \texttt{opencode run} como subprocessos reais; nenhuma simulação ou
\emph{mock} é utilizado.

\subsection{Corpus}
\label{sec:corpus}
Utilizamos uma amostra canônica do IMO-AnswerBench (R442, 4 problemas) +
cinco problemas auxiliares estilo-IMO com respostas determinísticamente
verificáveis (R503), totalizando $n=9$: algebra-001 (quociente de piso,
resposta 3), algebra-004 (desigualdade, $2^{u-2}$), teoria-dos-numeros-001
($p^3-q^5=(p+q)^2$, resposta $(7,3)$), combinatoria-001 (um máximo local,
$2^{n-1}$), algebra-002 ($2^n>n^2$, $n>1$, resposta 5), teoria-dos-numeros-002
($p^2 \mid 2^p+1$, $\{3\}$), teoria-dos-numeros-003 ($v_3(2023!)=1006$,
Legendre), combinatoria-002 (dois máximos locais em $S_5$, verificado 88
por enumeração independente), geometria-001 (diagonais de 2023-gono,
2043230). Os itens auxiliares são derivados estilo-IMO, não entradas
oficiais do Shortlist; identificamos accordingly para evitar erros de
atribuição.

\subsection{Resolvedor determinístico anti-vazamento}
\label{sec:solver}
O arsenal IMO original \citep{imo_harness_local} ecoava a resposta curta
esperada e auto-pontuava com nota 7 --- um vazamento. Implementamos o
\texttt{RealIMOSolver}, um resolvedor completamente determinístico (enumeração,
fórmula de Legendre, fórmula poligonal) cuja saída nunca contém a resposta
esperada antes do marcador ``Metodo:''; verificamos independência por testes
de ida-e-volta (nenhuma das 9 soluções vaza a resposta).

\subsection{Condições experimentais}
\label{sec:conditions}
Três condições no mesmo problema-prova NT-001: C0 (prompt bruto), C1
(apenas contexto), C2 (scaffold MASWOS de 16 estágios). A pontuação
composta é $S = 0,40A + 0,30V + 0,20C + 0,10D$, onde $A$=acurácia,
$V$=latência normalizada $1-(\ell-20)/90$ (limitada a $[0,1]$),
$C$=conformidade de formato, $D$=disponibilidade.

\subsection{Protocolo de validação cruzada (R503)}
\label{sec:crossval}
Após o ranking de problema único do R499, implementamos roteamento
automático (R500) e então re-executamos o experimento após reinício
completo do host (R503). O protocolo de validação cruzada é um delineamento
fatorial completo de 9 problemas $\times$ 3 modelos (sem divisão treino/
teste; todos os modelos veem todos os problemas, exatamente uma vez, em C2).
Reportamos:
\begin{itemize}
  \item acurácia por modelo $k/n$ com intervalos de confiança exatos
  Clopper--Pearson 95\% (binomial);
  \item testes McNemar pareados com correção de continuidade entre pares de
  modelos;
  \item $\kappa$ de Cohen (concordância além do acaso) para os dois modelos
  mais competentes;
  \item distribuições de latência por modelo (ok vs.\ timeout).
\end{itemize}

\subsection{Governança epistemológica (R504)}
\label{sec:governance}
A Camada de Governança Epistemológica adiciona quatro mecanismos de
integridade ao ecossistema:
\begin{enumerate}
  \item \textbf{Evidence Guard}: estados epistêmicos
  (OBSERVED/INFERRED/UNVERIFIED/BLOCKED); fontes aprendidas (policy,
  memória, confiança de skill) \emph{nunca} produzem OBSERVED.
  \item \textbf{Gates Feynman} (FEG-01..06): pontuação 0--12 baseada em
  mecanismo, evidência, separação observação/inferência, oracle refutável,
  necessidade demonstrada e experimento mínimo; FEG-07 (teach-back) é
  exclusivamente humano.
  \item \textbf{Avaliação Offline de Política v3}: decide$\rightarrow$observe
  separados (sem look-ahead), holdout temporal (70\%), Brier/ECE em
  shadow-matched-only, reward/regret com IC95\% bootstrap determinístico,
  detecção de drift e readiness sem auto-ativação.
  \item \textbf{Ledger de Auditoria SHA-256}: hash-chain encadeado,
  deduplicação por event\_id, verificação de integridade, snapshot somente
  leitura, exportação JSONL.
\end{enumerate}

\subsection{Orçamento de custo}
Cada chamada LLM é um subprocesso CLI real com timeout específico por modelo
(mimo 100\,s, big-pickle 110\,s, nemotron 150\,s). Todos os dados brutos,
tempos e status são armazenados incrementalmente (JSON) dentro do repositório
para auditoria.